\chapter{Grundlagen} \label{cha:fundamentals}

In diesem Teil der Arbeit werden die zum eindeutigen Verständnis unbedingt erforderlichen inhaltlichen Grundlagen und alle für das Thema relevanten Begriffsbestimmungen angeführt, um eine gemeinsame sprachliche und definitorische Basis zu schaffen.

Die Gliederungspunkte müssen möglichst prägnant bezeichnet werden.

Anmerkung: Die in dieser Formatvorlage beispielhaft enthaltenen (allgemein formulierten) Überschriften sind auf die im konkreten Fall tatsächlich passenden Überschriften und an die behandelten Themen anzupassen.

% --- Beispiele für Abbildung, Tabelle und Formel (können gelöscht werden) ---

\begin{figure}[htbp]
    \centering
    \includegraphics[width=0.4\textwidth]{00_figures/logo.jpg}
    \caption{Beispiel für eine Abbildung}
    \label{fig:example}
\end{figure}

\begin{table}[htbp]
    \centering
    \begin{tabular}{ll}
        \toprule
        {\fontsize{11pt}{12pt}\selectfont \textbf{Spalte 1}} & {\fontsize{11pt}{12pt}\selectfont \textbf{Spalte 2}} \\
        \midrule
        Wert A & Wert B \\
        Wert C & Wert D \\
        \bottomrule
    \end{tabular}
    \caption{Beispiel für eine Tabelle}
    \label{tab:example}
\end{table}

\begin{equation}
    U = R \cdot I
    \label{eq:ohm}
\end{equation}

Abbildung~\ref{fig:example}, Tabelle~\ref{tab:example} und Formel~\ref{eq:ohm} zeigen die Verwendung von Querverweisen.
\clearpage
